% Solutions of Book 4, chapter 18 (the network layer).
\section*{Chapter \ref{ch:b2:network-layer} --- The Network Layer: IP and Routing}

\begin{solution}{exo:b2:network-layer:1}
\code{198.51.100.77/27}: mask 255.255.255.224; $77 = 64 + 13$ and a \code{/27} block has 32 addresses, so the network
is 198.51.100.64, the broadcast 198.51.100.95, and $32 - 2 = 30$ hosts. \code{192.0.2.130/25}: mask 255.255.255.128,
network 192.0.2.128, broadcast 192.0.2.255, 126 hosts.
\end{solution}

\begin{solution}{exo:b2:network-layer:2}
The shift is by $32 - p + 1$ instead of $32 - p$, so the mask has $p - 1$ ones: for \code{/26} it is the mask of a
\code{/25}, 255.255.255.128, and the network comes out as 203.0.113.0. The membership test answers for that
\code{/25}: 203.0.113.10 is in 203.0.113.0/25, hence \code{True}, but it is not in 203.0.113.64/26. The answer is
wrong, and looks plausible because both \omterm{def:b2:network-layer:prefix}{prefixes} start alike. Testing against \code{ipaddress}, as \code{cidr.py}
does, catches it on the first address.
\end{solution}

\begin{solution}{exo:b2:network-layer:3}
\omcode{code/bachelor-2/18-network-layer/exercises/contains.py}{5}{8}{Membership in a \omterm{def:b2:network-layer:prefix}{prefix}}
Masking both sides makes the function correct even for a \omterm{def:b2:network-layer:prefix}{prefix} written with host bits set, such as
\code{203.0.113.77/26}. The test compares it with \code{ipaddress} on random addresses and lengths, 0 and 32
included.
\end{solution}

\begin{solution}{exo:b2:network-layer:4}
A \code{/26} holds $2^6 - 2 = 62$ machines. The offices need $s$ with $2^s - 2$ at least 50, 20, 10 and 2: $s = 6, 5,
4, 2$, that is a \code{/26}, a \code{/27}, a \code{/28} and a \code{/30}. Placing the largest first keeps each block
aligned on its size: 203.0.113.0/26, 203.0.113.64/27, 203.0.113.96/28 and 203.0.113.112/30, which leaves
203.0.113.116 to 203.0.113.255 free.
\end{solution}

\begin{solution}{exo:b2:network-layer:5}
\code{198.51.100.70} is in the \code{/24} and the \code{/26} (64 to 127), not in the \code{/27} (96 to 127): Y.
\code{198.51.100.100} is in all three: Z, the longest. \code{198.51.100.130} is only in the \code{/24}: X.
\code{192.0.2.1} matches only the default route: W.
\end{solution}

\begin{solution}{exo:b2:network-layer:6}
The new \omterm{def:b2:link-layer:edc}{checksum} is the complement of a sum that includes the old \omterm{def:b2:link-layer:edc}{checksum}, so it is not the complement of the sum of
the other words. The next \omterm{def:b2:network-layer:routing}{router} adds all ten words and gets \code{2b75} instead of \code{ffff}: it takes the header
for damaged and drops the packet. Every packet through this \omterm{def:b2:network-layer:routing}{router} is lost, which makes the bug obvious in practice,
and the output shows why: the receiver's sum is the check. Clearing the field first, as \code{forward} does, fixes it.
\end{solution}

\begin{solution}{exo:b2:network-layer:7}
\omcode{code/bachelor-2/18-network-layer/exercises/aggregate.py}{5}{23}{Aggregating \omterm{def:b2:network-layer:prefix}{prefixes}}
Two rules are applied until neither applies: a \omterm{def:b2:network-layer:prefix}{prefix} inside a shorter one of the set is dropped; two \omterm{def:b2:network-layer:prefix}{prefixes} of the
same length that differ only in their last bit, the two halves of a \omterm{def:b2:network-layer:prefix}{prefix} one bit shorter, are replaced by it. The
result covers the same addresses, and it is the fewest \omterm{def:b2:network-layer:prefix}{prefixes}: in the final set no \omterm{def:b2:network-layer:prefix}{prefix} is inside another and no
two are siblings, which is exactly the shape of the unique minimal cover. The test compares it with
\code{ipaddress.collapse_addresses} on random sets.
\end{solution}

\begin{solution}{exo:b2:network-layer:8}
Write $D_k(u, d)$ for the least cost of a path from $u$ to $d$ with at most $k$ links (infinite if there is none), and
assume that the value taken as infinity exceeds every distance. We show by induction on $k$ that after round $k$, the
table of $u$ gives $D_k(u, d)$ for every $d$, a destination being absent when the value is infinite.

For $k = 0$, the only path with no link is the empty path from $u$ to itself: $D_0(u, u) = 0$, and the initial table
holds exactly that. Suppose the claim holds after round $k$. A path from $u$ to $d \ne u$ with at most $k + 1$ links
starts with a link to some neighbour $v$, followed by a path from $v$ to $d$ with at most $k$ links; so
$D_{k+1}(u, d) = \min_v \bigl(c(u, v) + D_k(v, d)\bigr)$, and $D_{k+1}(u, u) = 0$ since costs are positive. Round $k +
1$ computes exactly this minimum from the neighbours' tables of round $k$, which by hypothesis hold $D_k(v, d)$.

With positive costs a shortest path visits no \omterm{def:b2:network-layer:routing}{router} twice, so it has at most $n - 1$ links: $D_{n-1}(u, d)$ is the
shortest distance, and $D_k = D_{n-1}$ for every $k \ge n - 1$. The distances therefore stop changing after at most
$n - 1$ rounds; the next hops, chosen by the same rule from the same values, stop with them.
\end{solution}

\begin{solution}{exo:b2:network-layer:9}
\omcode{code/bachelor-2/18-network-layer/listings/fragment.py}{4}{12}{\omterm{def:b2:network-layer:fragmentation}{Fragmentation}}
With an \omterm{def:b2:network-layer:fragmentation}{MTU} of 1006, each fragment carries at most $(1006 - 20)$ rounded down to a multiple of 8, that is 984 bytes:
fragments of 984, 984, 984 and 48 bytes, at offsets 0, 123, 246 and 369 (in units of 8 bytes), the flag MF set on all
but the last. The offset field has 13 bits and counts 8-byte units, so offsets up to $(2^{13} - 1) \times 8 =
65\,528$ bytes cover any packet, whose total length is at most $65\,535$ bytes; this saves 3 bits of header, and requires every fragment but the last to end on a multiple of 8.
\end{solution}

\begin{solution}{exo:b2:network-layer:10}
\omoutput{code/bachelor-2/18-network-layer/outputs/poison.txt}{Plain and poisoned, after the
failure.}
The change is one line in \code{dv_round}: when $v$'s route to $d$ goes through $u$, $u$ reads it as infinite
(\code{poison=True}). On the line, A routes to C through B, so it tells B ``infinity''; B, having lost its link, has
no route left and says so at once, and both reach 16 in 2 rounds. In the triangle, after C--D fails, A and B both
route to D through C, so each poisons only its announcement to C, and C has no route left. But A still tells B ``D
at 2'': A's route goes through C, not through B, so nothing is poisoned. In round 2, A and B each take D at 3 through
the other. In round 3 they poison each other and drop to 16, but C hears A's 3, which does not go through C, and
takes D at 4 through A; in round 4, B takes 5 through C. The stale route travels round the triangle, one hop and one
unit per round. Poisoning only stops a \omterm{def:b2:network-layer:routing}{router} from echoing a route back to its own next hop, a loop of two; this loop
has three, and the \omterm{def:b2:network-layer:routing}{routers} count to infinity in 15 rounds, as without it.
\end{solution}

\begin{solution}{exo:b2:network-layer:11}
A linear table costs one test per \omterm{def:b2:network-layer:prefix}{prefix}: $n$ tests, as the printed times show (about 0.9 µs per \omterm{def:b2:network-layer:prefix}{prefix}, four times
the time for four times the \omterm{def:b2:network-layer:prefix}{prefixes}). A trie costs one step per bit of the longest matching path, at most 32,
whatever $n$; its time rises a little with $n$ only because more \omterm{def:b2:network-layer:prefix}{prefixes} make longer paths. Binary search does not
work directly on nested \omterm{def:b2:network-layer:prefix}{prefixes}, since the longest match need not be next to the address in sorted order. Each
\omterm{def:b2:network-layer:prefix}{prefix} is a range of addresses; the $n$ ranges cut the address space into at most $2n + 1$ intervals, each labelled
with the longest \omterm{def:b2:network-layer:prefix}{prefix} that contains it. A binary search among these intervals costs $\log_2(2n + 1)$ steps, about
13 for 4000 \omterm{def:b2:network-layer:prefix}{prefixes}, but an update may relabel many intervals, while a trie changes one node.
\end{solution}

\begin{solution}{exo:b2:network-layer:12}
\omcode{code/bachelor-2/18-network-layer/exercises/ipv6_compress.py}{4}{22}{The recommended text form}
The loop finds the longest run of zero groups, taking a run only if it is strictly longer than the best so far, so the
first of two equal runs wins and a run of one group (the initial \code{best_len} is 1) is never shortened. The
groups are written in lowercase hexadecimal without leading zeros. The test draws addresses whose groups are often 0,
so that ties and runs at either end occur, and compares with \code{str(ipaddress.IPv6Address(...))}.
\end{solution}

\begin{solution}{pb:b2:network-layer:1}
\textbf{1.} $256 - 30 = 226$ addresses, two of which (\code{.0} and \code{.255}) are the network and broadcast
addresses if the whole block is one network.

\textbf{2.} Five machines plus the network and broadcast addresses need 7 addresses: a \code{/29} (8 addresses, 6
machines). A \code{/24} holds $2^{29-24} = 32$ of them.

\textbf{3.} No. The \omterm{def:b2:network-layer:routing}{router} addresses 1 to 30 fall in the first four \code{/29}s (0 to 31), which leaves 28 for 30
networks.

\textbf{4.} Place $k$ gets $8k$ to $8k + 7$; for $k = 30$ that is 240 to 247, inside the block, and the 30 ranges do
not overlap. Place 12: \code{203.0.113.96/29}.

\textbf{5.} Network 203.0.113.96, broadcast 203.0.113.103.

\textbf{6.} The place number is the last byte of an address divided by 8: an address tells at once where it is.

\textbf{7.} The Library's \omterm{def:b2:network-layer:routing}{router}, \code{198.51.100.3}, the first hop.

\textbf{8.} 30 entries, one per place, plus a default route. Seen from outside, the whole town is one entry,
\code{203.0.113.0/24}: aggregation hides the 30 \omterm{def:b2:network-layer:prefix}{prefixes}.

\textbf{9.} \code{203.0.113.0/24} (and the \omterm{def:b2:network-layer:routing}{routers}' \code{198.51.100.0/24}) towards the inside, and \code{0.0.0.0/0}
towards the provider.

\textbf{10.} $97 = 8 \times 12 + 1$: place 12.

\textbf{11.} 29 steps: the walk follows the address's bits to the node of \code{203.0.113.96/29} at depth 29, which
has no children.

\textbf{12.} The \omterm{def:b2:network-layer:ttl}{time to live} must exceed the length of any path a packet may need, anywhere on the Internet, not just
in town; 64 leaves a wide margin while still killing a looping packet quickly.

\textbf{13.} At most $n - 1 = 29$ rounds. The tables stop changing once every shortest path is known, so 8 is the
largest number of links on a shortest path (choosing, for each pair, the shortest path with fewest links).

\textbf{14.} With hop counts the shortest paths are those with fewest links, and the longest of them has 7 links; with
lengths, the shortest path may take more, shorter, links.

\textbf{15.} Each \omterm{def:b2:network-layer:routing}{router} sends its table to each neighbour: $2 \times 53 = 106$ messages per round.

\textbf{16.} The whole map: the 53 links (each described by both of its \omterm{def:b2:network-layer:routing}{routers}, 106 descriptions).

\textbf{17.} With a heap, about $(n + m)\log_2 n \approx (30 + 106) \times 5 \approx 700$ elementary steps: nothing for
a \omterm{def:b2:network-layer:routing}{router}.

\textbf{18.} Link state: the map is small, each \omterm{def:b2:network-layer:routing}{router} computes its routes in an instant, and a failure costs one
flood and one recomputation instead of a count to infinity.

\textbf{19.} Distance vector has no message that says ``this \omterm{def:b2:network-layer:routing}{router} is gone''. The café's neighbours lose their links
to it, but still hear the other neighbours announce routes to it, routes which went through the café's own link and
are now stale.

\textbf{20.} The nearest \omterm{def:b2:network-layer:routing}{routers} were at 1, directly linked. After the failure their best offer comes from another
former neighbour, which announced 1 in the previous round: 2. In the next round every offer is at least 2, so the
least distance is 3, and so on: the least distance grows by one link per round, since every route now passes through
another \omterm{def:b2:network-layer:routing}{router} that believes in an older, shorter one.

\textbf{21.} After round $r$, the least distance is $r + 1$ (by the argument of question 20), so all reach 16 when $r
+ 1 = 16$, at round 15, whatever the distances were before: a far \omterm{def:b2:network-layer:routing}{router} keeps its distance until the rising minimum $r + 1$
overtakes it, then rises with the rest, and all reach 16 together.

\textbf{22.} $15 \times 30$~s, seven and a half minutes, during which packets for the café circle until their \omterm{def:b2:network-layer:ttl}{time to
live} runs out.

\textbf{23.} The café's neighbours flood the loss of their links; each \omterm{def:b2:network-layer:routing}{router} removes them from its map and runs
Dijkstra once, finding the café unreachable. It takes the time of one flood, a fraction of a second.

\textbf{24.} Count to infinity: 15 rounds to reach the limit 16 on the town's network, against one recomputation per
\omterm{def:b2:network-layer:routing}{router} in link state.
\end{solution}
